% Source: https://artofproblemsolving.com/community/c3h2111553_continued_fraction_equals_continued_fraction
% Stable identity: aops-topic-2111553-q01
\textbf{A shortcut.} Let $y=7+\dfrac{1}{7+\dfrac{1}{7}}$, so the left side is $1-\dfrac{1}{y}$. For any $y>1$,
\[1-\dfrac{1}{y}=\dfrac{y-1}{y}=\dfrac{1}{\dfrac{y}{y-1}}=\dfrac{1}{1+\dfrac{1}{y-1}},\qquad y-1=6+\dfrac{1}{7+\dfrac{1}{7}}.\]
So the left side equals
\[\dfrac{1}{1+\dfrac{1}{6+\dfrac{1}{7+\dfrac{1}{7}}}},\]
which has the required form with $a=1$, $b=6$, $c=7$, $d=7$.

\textbf{Uniqueness.} Numerically, $y=7+\dfrac{7}{50}=\dfrac{357}{50}$ and the left side is $\dfrac{307}{357}$. For positive integers, $c+\dfrac{1}{d}>1$, so $b<b+\dfrac{1}{c+\dfrac{1}{d}}<b+1$; similarly $a$ is the integer part of $a+\dfrac{1}{b+\cdots}=\dfrac{357}{307}$, so $a=1$. Then $b+\dfrac{1}{c+\dfrac{1}{d}}=\dfrac{307}{50}$ gives $b=6$, and $c+\dfrac{1}{d}=\dfrac{50}{7}$. If $d=1$, then $c+1=\dfrac{50}{7}$ is not an integer; so $d\ge2$, $c$ is the integer part $7$, and $\dfrac{1}{d}=\dfrac{1}{7}$, $d=7$. Hence the values are forced, and
\[abcd=1\cdot6\cdot7\cdot7=\boxed{294}.\]
